\begin{insight}{逆矩阵是把变形倒着做}
矩阵的每一列是标准基向量变换后的去向。对$S=\begin{pmatrix}1&1\\0&1\end{pmatrix}$，竖直单位边变成$(1,1)\T$，水平边不动，单位正方形被剪成平行四边形。逆变换把每个点的横坐标减去其纵坐标，因此$S\inv=\begin{pmatrix}1&-1\\0&1\end{pmatrix}$。
\end{insight}
\begin{figure}[H]\centering
\begin{tikzpicture}[scale=1.35]
\begin{scope}
\draw[laaxis](-.2,0)--(2.4,0);\draw[laaxis](0,-.2)--(0,1.6);
\filldraw[fill=Teal!12,draw=Teal](0,0)rectangle(1,1);
\draw[lavec,Ink](0,0)--(1,0)node[midway,below]{$e_1$};\draw[lavec,Gold](0,0)--(0,1)node[left]{$e_2$};
\node[lanote]at(.7,-.5){原坐标};
\end{scope}
\draw[-{Stealth},Teal,thick](2.5,.75)--(3.8,.75)node[midway,above]{$S$};
\draw[-{Stealth},Gold,thick](3.8,.25)--(2.5,.25)node[midway,below]{$S\inv$};
\begin{scope}[xshift=4.3cm]
\draw[laaxis](-.2,0)--(2.4,0);\draw[laaxis](0,-.2)--(0,1.6);
\filldraw[fill=Teal!12,draw=Teal](0,0)--(1,0)--(2,1)--(1,1)--cycle;
\draw[lavec,Ink](0,0)--(1,0)node[midway,below]{$Se_1$};\draw[lavec,Gold](0,0)--(1,1)node[above]{$Se_2$};
\node[lanote]at(.9,-.5){剪切后的坐标};
\end{scope}
\end{tikzpicture}
\caption{先作$S$再作$S\inv$，每个向量回到原处，而不只是面积恢复。}\label{fig:matrix-inverse}
\end{figure}
\begin{clevertrick}{用幂零看穿剪切矩阵}
写$S=I+N$，其中$N=\begin{pmatrix}0&1\\0&0\end{pmatrix}$，直接相乘得$N^2=0$。当$m\ge0$时，$I,N$可交换，二项式展开中的二次及以上$N$项全为零，故
\[S^m=\sum_{j=0}^m\binom mjN^j=I+mN=\begin{pmatrix}1&m\\0&1\end{pmatrix}.\]
又$(I+N)(I-N)=I-N^2=I$，反序相乘也为$I$，所以$S\inv=I-N$。若$m=-k<0$，则$S^m=(I-N)^k=I-kN=I+mN$，幂公式仍成立。例如
\[S^{100}(0,1)\T=(0,1)\T+100N(0,1)\T=(0,1)\T+100(1,0)\T=(100,1)\T.\]
一般块矩阵$S=I+N$若只有一个非零副对角块，先核验$N^2=0$，便可逐字使用同一推导。
\textbf{易错提示：}面积不变不表示矩阵等于单位阵；图中剪切的行列式为1，但大部分向量都改变了方向。
\end{clevertrick}
